\section{Splitting by the sign of the symbol}\label{sec:compression}
The companion manuscript presents the target as a compression, $Q_L=K_+-T_L$,
with $K_+$ the archimedean energy together with the positive part of the
constant $c_L$ of \eqref{eq:sigma} and the even half of the pole form, and $T_L$
the negative part of that constant, the odd half of the pole form and the mirror
prime references; both sides are positive term by term. That presentation moves
the whole constant to whichever side its sign puts it, at every frequency. By
Section~\ref{sec:density} only the band $|\tau|<\tau_L$ has to move.

\begin{proposition}[Symbol split]\label{prop:symbol-split}
Let $\sigma_L^\pm=\max(\pm\sigma_L,0)$ for $\sigma_L$ as in \eqref{eq:sigma},
and define on $C_c^\infty(I_L)$
\begin{align}
 K_+^{\rm new}[f]&=\frac1{2\pi}\int_\R\sigma_L^+(\tau)|\widehat{E_Lf}|^2\dd\tau
  +2\Big|\int_{I_L}f\cosh(x/2)\Big|^2,\label{eq:source-new}\\
 T_L^{\rm new}[f]&=\frac1{2\pi}\int_\R\sigma_L^-(\tau)|\widehat{E_Lf}|^2\dd\tau
  +2\Big|\int_{I_L}f\sinh(x/2)\Big|^2+\sum_{p\in\mathcal P}\widetilde B_{r_p,d_p}[E_Lf].
 \label{eq:subtracted-new}
\end{align}
Then $Q_L=K_+^{\rm new}-T_L^{\rm new}$ identically, both sides are positive term
by term, and $K_+^{\rm new}\preceq K_+$, strictly except at $\tau=0$.
\end{proposition}
\begin{proof}
Start from the prime-free subtraction form $Q_L=A_L-\sum_p\widetilde B_{r_p,d_p}$
of the companion manuscript, in which $A_L$ is the archimedean energy together
with the constant $c_L$ and the pole form. By \eqref{eq:sigma} its first two
terms are the Fourier multiplier with symbol $\sigma_L$; split that as
$\sigma_L^+-\sigma_L^-$ and split $P_L$ by \eqref{eq:parity-split}. A multiplier
with nonnegative symbol is a positive form, and each $\widetilde B_{r_p,d_p}$ is
positive. For the comparison, $\sigma_L^-$ attains $(c_L)_-$ only at $\tau=0$
because $2\theta'$ is strictly increasing on $\tau>0$, and $K_+-K_+^{\rm new}$ is
the multiplier with symbol $(c_L)_--\sigma_L^-\geq0$.
\end{proof}

The earlier presentation is the case in which $\sigma_L^-$ is replaced by the
constant $(c_L)_-$ everywhere. In a $48$-cell basis the least generalized
eigenvalue improves as follows, both presentations assembled independently and
both identities verified to machine zero.
\begin{center}
\small
\begin{tabular}{lrrrrrrr}
\toprule
$L$ & $4/5$ & $1$ & $5/4$ & $3/2$ & $2$ & $5/2$ & $3$\\
\midrule
$\min\operatorname{spec}(Q_L;K_+)$
 & $3.68$e$-4$ & $3.68$e$-4$ & $2.81$e$-4$ & $2.65$e$-4$ & $1.81$e$-4$ & $1.89$e$-4$ & $1.15$e$-4$\\
$\min\operatorname{spec}(Q_L;K_+^{\rm new})$
 & $1.45$e$-3$ & $1.25$e$-3$ & $6.84$e$-4$ & $5.68$e$-4$ & $2.45$e$-4$ & $2.11$e$-4$ & $1.15$e$-4$\\
ratio & $3.93$ & $3.40$ & $2.43$ & $2.15$ & $1.35$ & $1.11$ & $1.00$\\
\bottomrule
\end{tabular}
\end{center}

These are one-sided subspace values: above one they would certify that the
domination fails, below one they certify nothing. \emph{The gain is a constant
factor and not a change of kind}: the domination remains saturated and the
contraction remains critical. What the split buys is structural. The constant
$w_0$ never appears as an independent object; for $L>\log13$ the subtracted side
is only the odd pole term and the mirror references, with no constant at all;
and the band that has to move is explicit and shrinking.
